% Options for packages loaded elsewhere
\PassOptionsToPackage{unicode}{hyperref}
\PassOptionsToPackage{hyphens}{url}
\documentclass[
]{book}
\usepackage{xcolor}
\usepackage{amsmath,amssymb}
\setcounter{secnumdepth}{5}
\usepackage{iftex}
\ifPDFTeX
  \usepackage[T1]{fontenc}
  \usepackage[utf8]{inputenc}
  \usepackage{textcomp} % provide euro and other symbols
\else % if luatex or xetex
  \usepackage{unicode-math} % this also loads fontspec
  \defaultfontfeatures{Scale=MatchLowercase}
  \defaultfontfeatures[\rmfamily]{Ligatures=TeX,Scale=1}
\fi
\usepackage{lmodern}
\ifPDFTeX\else
  % xetex/luatex font selection
\fi
% Use upquote if available, for straight quotes in verbatim environments
\IfFileExists{upquote.sty}{\usepackage{upquote}}{}
\IfFileExists{microtype.sty}{% use microtype if available
  \usepackage[]{microtype}
  \UseMicrotypeSet[protrusion]{basicmath} % disable protrusion for tt fonts
}{}
\makeatletter
\@ifundefined{KOMAClassName}{% if non-KOMA class
  \IfFileExists{parskip.sty}{%
    \usepackage{parskip}
  }{% else
    \setlength{\parindent}{0pt}
    \setlength{\parskip}{6pt plus 2pt minus 1pt}}
}{% if KOMA class
  \KOMAoptions{parskip=half}}
\makeatother
\usepackage{longtable,booktabs,array}
\usepackage{calc} % for calculating minipage widths
% Correct order of tables after \paragraph or \subparagraph
\usepackage{etoolbox}
\makeatletter
\patchcmd\longtable{\par}{\if@noskipsec\mbox{}\fi\par}{}{}
\makeatother
% Allow footnotes in longtable head/foot
\IfFileExists{footnotehyper.sty}{\usepackage{footnotehyper}}{\usepackage{footnote}}
\makesavenoteenv{longtable}
\usepackage{graphicx}
\makeatletter
\newsavebox\pandoc@box
\newcommand*\pandocbounded[1]{% scales image to fit in text height/width
  \sbox\pandoc@box{#1}%
  \Gscale@div\@tempa{\textheight}{\dimexpr\ht\pandoc@box+\dp\pandoc@box\relax}%
  \Gscale@div\@tempb{\linewidth}{\wd\pandoc@box}%
  \ifdim\@tempb\p@<\@tempa\p@\let\@tempa\@tempb\fi% select the smaller of both
  \ifdim\@tempa\p@<\p@\scalebox{\@tempa}{\usebox\pandoc@box}%
  \else\usebox{\pandoc@box}%
  \fi%
}
% Set default figure placement to htbp
\def\fps@figure{htbp}
\makeatother
\setlength{\emergencystretch}{3em} % prevent overfull lines
\providecommand{\tightlist}{%
  \setlength{\itemsep}{0pt}\setlength{\parskip}{0pt}}
\usepackage{bookmark}
\IfFileExists{xurl.sty}{\usepackage{xurl}}{} % add URL line breaks if available
\urlstyle{same}
\hypersetup{
  pdftitle={Seminar Notes 2026},
  pdfauthor={Yi Zhou},
  hidelinks,
  pdfcreator={LaTeX via pandoc}}

\title{Seminar Notes 2026}
\author{Yi Zhou}
\date{2026-04-14}

\begin{document}
\maketitle

{
\setcounter{tocdepth}{1}
\tableofcontents
}
\chapter*{資料}\label{ux8cc7ux6599}
\addcontentsline{toc}{chapter}{資料}

\section*{スケジュール}\label{ux30b9ux30b1ux30b8ux30e5ux30fcux30eb}
\addcontentsline{toc}{section}{スケジュール}

\href{Schedule2026.xlsx}{Click here to download XLSX}

\section*{参考資料}\label{ux53c2ux8003ux8cc7ux6599}
\addcontentsline{toc}{section}{参考資料}

\begin{itemize}
\tightlist
\item
  Statistical Inference \href{Textbook/Casella_Berger_Statistical_Inference.pdf}{Click here to open PDF}

  \begin{itemize}
  \tightlist
  \item
    現代数理統計学 \href{https://elib.maruzen.co.jp/elib/html/BookDetail/Id/3000105174;jsessionid=77611339D268556CDFB560D03E30F354?0}{Link to Kobe Univ. Lib}
  \item
    現代数理統計学の基礎 \href{https://elib.maruzen.co.jp/elib/html/BookDetail/Id/3000050917?2}{Link to Kobe Univ. Lib}
  \item
    一般化線形モデル \href{https://elib.maruzen.co.jp/elib/html/BookDetail/Id/3000031957?19}{Link to Kobe Univ. Lib}
  \item
    ベイズ計算統計学 \href{https://elib.maruzen.co.jp/elib/html/BookDetail/Id/3000027079?21}{Link to Kobe Univ. Lib}
  \end{itemize}
\item
  Fundamentals of Biostatistics \href{Textbook/Fundamentals\%20of\%20Biostatistics\%20(8th\%20Edition).pdf}{Click here to open PDF}

  \begin{itemize}
  \tightlist
  \item
    医学統計学ハンドブック \href{https://elib.maruzen.co.jp/elib/html/BookDetail/Id/3000068387?3}{Link to Kobe Univ. Lib}
  \item
    生存時間解析 \href{https://elib.maruzen.co.jp/elib/html/BookDetail/Id/3000117213?16}{Link to Kobe Univ. Lib}
  \item
    実験計画法と分散分析 \href{https://elib.maruzen.co.jp/elib/html/BookDetail/Id/3000027080?23}{Link to Kobe Univ. Lib}
  \end{itemize}
\item
  An introduction to statistical learning

  \begin{itemize}
  \tightlist
  \item
    with application in Python \href{Textbook/ISLP_pyhton.pdf}{Click here to open PDF}
  \item
    with application in R \href{Textbook/ISLR_R.pdf}{Click here to open PDF}
  \end{itemize}
\item
  統計的因果推論 \href{https://elib.maruzen.co.jp/elib/html/BookDetail/Id/3000034381?17}{Link to Kobe Univ. Lib}

  \begin{itemize}
  \tightlist
  \item
    欠測データの統計解析 \href{https://elib.maruzen.co.jp/elib/html/BookDetail/Id/3000031956?18}{Link to Kobe Univ. Lib}
  \end{itemize}
\item
  Others \href{Textbook/.}{Click here to open folder}
\end{itemize}

\section*{数理科学入門Slides}\label{ux6570ux7406ux79d1ux5b66ux5165ux9580slides}
\addcontentsline{toc}{section}{数理科学入門Slides}

\begin{itemize}
\tightlist
\item
  Slides \href{Stats-slides/allslides.pdf}{Click here to open PDF}
\item
  Notes \href{Stats-slides/allnotes.pdf}{Click here to open PDF}
\item
  Others \href{Stats-slides/.}{Click here to open folder}
\end{itemize}

\section*{R}\label{r}
\addcontentsline{toc}{section}{R}

\begin{itemize}
\tightlist
\item
  R and RStudio installation: \url{https://posit.co/download/rstudio-desktop/}
\item
  授業の資料：\url{https://lec25.github.io/2026-1H609/}
\item
  etc.
\end{itemize}

\section*{Python}\label{python}
\addcontentsline{toc}{section}{Python}

\begin{itemize}
\tightlist
\item
  \url{https://yongks.github.io/python_bookdown/}
\item
  etc.
\end{itemize}

\section*{Journals}\label{journals}
\addcontentsline{toc}{section}{Journals}

\begin{itemize}
\item
  Biostatistics

  \begin{itemize}
  \tightlist
  \item
    \href{https://academic.oup.com/biometrics}{Biometrics}
  \item
    \href{https://onlinelibrary.wiley.com/journal/10970258}{Statistics in Medicine}
  \item
    \href{https://academic.oup.com/biostatistics}{Biostatistics}
  \item
    \href{https://journals.sagepub.com/home/SMM}{Statistical Methods in Medical Research (SMMR)}
  \item
    \href{https://www.tandfonline.com/journals/uasa20}{Journal of the American Statistical Association (JASA)}
  \item
    \href{https://academic.oup.com/biomet}{Biometrika}
  \item
    \href{https://academic.oup.com/jrsssb?login=true}{Journal of the Royal Statistical Society, Series B (Statistical Methodology) (JRSS B)}
  \item
    \href{https://www.jstage.jst.go.jp/browse/jjb/-char/ja}{計量生物学}
  \item
    \href{https://www.ism.ac.jp/editsec/toukei/index.html}{統計数理}
  \end{itemize}
\item
  Epidemiology

  \begin{itemize}
  \tightlist
  \item
    \href{https://academic.oup.com/aje}{American Journal of Epidemiology}
  \end{itemize}
\item
  Bioinformatics

  \begin{itemize}
  \tightlist
  \item
    \href{https://academic.oup.com/bib}{Briefings in Bioinformatics}
  \item
    \href{https://academic.oup.com/bioinformatics}{Bioinformatics}
  \end{itemize}
\end{itemize}

\chapter{Exponential family（指数型分布族）}\label{exponential-familyux6307ux6570ux578bux5206ux5e03ux65cf}

2026-4-15

\section*{復習}\label{ux5fa9ux7fd2}
\addcontentsline{toc}{section}{復習}

\begin{itemize}
\tightlist
\item
  確率空間、確率変数
\item
  離散型、連続型確率変数 (Random Variable, RV)
\item
  確率(密度)関数、分布：Probaility Mass Function (PMF) vs.~Probability Density Function (PDF) \(\rightarrow\) Cumulative Distribution Function (CDF)
\item
  離散型、連続型代表的な分布、分布の形、近似関係
\item
  分布の期待値、分散 \(\rightarrow\) Moment (Moment Generating Function)
\item
  Random variable vs.~Random sample（無作為標本）
\item
  PDF, PMF vs.~Likelihood（尤度関数）
\end{itemize}

\section{指数型分布族（Statistical Inference pp.111\textasciitilde；現代数理統計の基礎 pp.119）}\label{ux6307ux6570ux578bux5206ux5e03ux65cfstatistical-inference-pp.111ux73feux4ee3ux6570ux7406ux7d71ux8a08ux306eux57faux790e-pp.119}

\begin{itemize}
\tightlist
\item
  数式１
\item
  連続型の例：Normal（証明）, Gamma, Beta, Exponential, Chi-square
\item
  離散型の例：Binomial（証明）, Bernoulli, Poisson, Negative binomial, Multinomial
\end{itemize}

\subsection{定理１}\label{ux5b9aux7406uxff11}

\begin{itemize}
\tightlist
\item
  期待値、分散の導出（課題：証明 \(\rightarrow\) Exercise 3.31, 3.32）

  \begin{itemize}
  \tightlist
  \item
    例：Binomial, Normal exponential family （証明）
  \end{itemize}
\end{itemize}

\subsection{数式２}\label{ux6570ux5f0fuxff12}

\begin{itemize}
\tightlist
\item
  Normal exponential family
\item
  Curved exponential family*

  \begin{itemize}
  \tightlist
  \item
    Poisson distribution（例）*
  \end{itemize}
\end{itemize}

\subsection{指数型分布族 Random Sampleの和（pp.217）}\label{ux6307ux6570ux578bux5206ux5e03ux65cf-random-sampleux306eux548cpp.217}

\begin{itemize}
\tightlist
\item
  数式
\item
  例： Sum of Bernoulli RVs
\end{itemize}

\section{Location and Scale family（pp.116\textasciitilde）}\label{location-and-scale-familypp.116}

\begin{itemize}
\tightlist
\item
  定理： Location, scaleによってPDFのを構成する
\end{itemize}

\subsection{Location familyの定義}\label{location-familyux306eux5b9aux7fa9}

\begin{itemize}
\tightlist
\item
  例： Exponential location family
\end{itemize}

\subsection{Scale family, Location-Scale family}\label{scale-family-location-scale-family}

\begin{itemize}
\item
  定理１：分布の導出

  \(X=\sigma Z+\mu \Rightarrow f_Z(z)\rightarrow f_X(x)\)
\item
  定理２：\(X\) の期待値、分散

  \begin{itemize}
  \tightlist
  \item
    性質：\(X\) の確率
  \end{itemize}
\item
  応用：Location-Scale pivot (pp.427；現 pp.172)　for constructing CI
\end{itemize}

\section*{コメント}\label{ux30b3ux30e1ux30f3ux30c8}
\addcontentsline{toc}{section}{コメント}

\begin{itemize}
\tightlist
\item
  確率分布を統一的に表現
\item
  共通の性質を解析可能
\item
  期待値・分散の導出が容易
\item
  データとパラメータの分離 --\textgreater{} Sufficient Statistics（十分統計量、後に）
\end{itemize}

\section*{応用}\label{ux5fdcux7528}
\addcontentsline{toc}{section}{応用}

\begin{itemize}
\tightlist
\item
  Generalized linear model（一般化線形モデル）

  \begin{itemize}
  \tightlist
  \item
    Exponential familyの標準形（canonical form） --\textgreater{} link function
  \end{itemize}
\end{itemize}

\chapter{Hierarchical Model, Mixture Distribution（階層モデル、混合分布）}\label{hierarchical-model-mixture-distributionux968eux5c64ux30e2ux30c7ux30ebux6df7ux5408ux5206ux5e03}

2026-4-22

\section{復習}\label{ux5fa9ux7fd2-1}

\begin{itemize}
\tightlist
\item
  条件付き確率（密度）関数、周辺確率（密度）関数
\item
  Bayesian Rule, Bayes estimators
\end{itemize}

\section{階層モデル（pp.163\textasciitilde）}\label{ux968eux5c64ux30e2ux30c7ux30ebpp.163}

\begin{itemize}
\item
  例： Binomial-Poisson Hierarchy

  \begin{itemize}
  \tightlist
  \item
    分布と期待値
  \end{itemize}
\end{itemize}

\subsection{定理１：Conditional expectation（証明）}\label{ux5b9aux7406uxff11conditional-expectationux8a3cux660e}

\section{Mixture Distributionの定義}\label{mixture-distributionux306eux5b9aux7fa9}

\begin{itemize}
\item
  例： Normal scale mixture distribution (現 pp.65)
\item
  例： Binomial-Poisson Hierarchy

  \begin{itemize}
  \tightlist
  \item
    分布と期待値
  \end{itemize}
\item
  例： Noncentral Chi2 distribution
\item
  例： Beta-binomial Hierarchy
\end{itemize}

\subsection{定理２：Conditional variance identity（証明）}\label{ux5b9aux7406uxff12conditional-variance-identityux8a3cux660e}

\begin{itemize}
\tightlist
\item
  例： Beta-binomial Hierarchy
\end{itemize}

\subsection{EM algorithm (pp.327)}\label{em-algorithm-pp.327}

\begin{itemize}
\tightlist
\item
  例： Multiple Poisson rates
\item
  定理*
\item
  EM vs MLE*
\item
  Missing data issue*
\end{itemize}

\section*{コメント}\label{ux30b3ux30e1ux30f3ux30c8-1}
\addcontentsline{toc}{section}{コメント}

\begin{itemize}
\tightlist
\item
  Posterior 分布の導出
\item
  不確実性を階層的にモデル化
\end{itemize}

\section*{応用}\label{ux5fdcux7528-1}
\addcontentsline{toc}{section}{応用}

\begin{itemize}
\item
  Bayes analysis (pp.371)

  \begin{itemize}
  \tightlist
  \item
    Empirical Bayes analysis
  \item
    Hierarchical Bayes analysis
  \item
    Gibbs sampling
  \end{itemize}
\item
  GLMM: generalized linear mixed effects model（一般化線形混合モデル）

  \begin{itemize}
  \tightlist
  \item
    その特例： ANOVA
  \item
    その一般化： Bayesian hierarchical model、階層ベイズ
  \end{itemize}
\item
  Meta-analysis　データの統合
\end{itemize}

\chapter{Order Statistics（順序統計量）}\label{order-statisticsux9806ux5e8fux7d71ux8a08ux91cf}

\section*{復習}\label{ux5fa9ux7fd2-2}
\addcontentsline{toc}{section}{復習}

\begin{itemize}
\tightlist
\item
  統計量
\item
  Rank
\item
  Nonparametric test (based on rank)

  \begin{itemize}
  \tightlist
  \item
    Sign test/Wilcoxon Signed-rank test
  \item
    Wilcoxon rank-sum test/U test
  \end{itemize}
\end{itemize}

\section{順序統計量（pp.226\textasciitilde；現 pp.101）}\label{ux9806ux5e8fux7d71ux8a08ux91cfpp.226ux73fe-pp.101}

\begin{itemize}
\item
  定義、表記

  \begin{itemize}
  \tightlist
  \item
    Sample range, sample median
  \item
    lower/upper quartile, interquartile range
  \end{itemize}
\end{itemize}

\section{定理1：順序統計量の確率（導出）}\label{ux5b9aux74061ux9806ux5e8fux7d71ux8a08ux91cfux306eux78baux7387ux5c0eux51fa}

\section{定理2：順序統計量のCDF, PDF（導出）}\label{ux5b9aux74062ux9806ux5e8fux7d71ux8a08ux91cfux306ecdf-pdfux5c0eux51fa}

\begin{itemize}
\tightlist
\item
  例： Uniform distribution
\end{itemize}

\section{定理3：順序統計量のJoint CDF, PDF（Exercise 5.27, 5.28）}\label{ux5b9aux74063ux9806ux5e8fux7d71ux8a08ux91cfux306ejoint-cdf-pdfexercise-5.27-5.28}

\begin{itemize}
\tightlist
\item
  例： Midrange, range
\end{itemize}

\section*{応用}\label{ux5fdcux7528-2}
\addcontentsline{toc}{section}{応用}

\begin{itemize}
\item
  Kolmogorov--Smirnov 検定 for Goodness of Fit（適合度・あてはまりの良さ)

  \begin{itemize}
  \tightlist
  \item
    Goodness-of-Fit Tests: Chi2 test
  \end{itemize}
\end{itemize}

\chapter{Generating a random sample（乱数の生成）}\label{generating-a-random-sampleux4e71ux6570ux306eux751fux6210}

\section*{復習}\label{ux5fa9ux7fd2-3}
\addcontentsline{toc}{section}{復習}

\begin{itemize}
\tightlist
\item
  Law of large numbers (LLN)
\item
  Convergence concept:

  \begin{itemize}
  \tightlist
  \item
    in prob. (weak LLN)
  \item
    almost surely (strong LLN)
  \item
    in distribution (central limit theorem, CLT)
  \end{itemize}
\item
  変数変換
\end{itemize}

\section{例： Exponential lifetime (pp.245\textasciitilde)}\label{ux4f8b-exponential-lifetime-pp.245}

\section{直接法}\label{ux76f4ux63a5ux6cd5}

\begin{itemize}
\tightlist
\item
  Prob integral transformation
\item
  Box-Muller algorithm
\item
  Poisson variance
\end{itemize}

\section{間接法}\label{ux9593ux63a5ux6cd5}

\begin{itemize}
\tightlist
\item
  Beta RV
\end{itemize}

\section{Accept/Reject Algorithm}\label{acceptreject-algorithm}

\begin{itemize}
\item
  Beta RV
\item
  Metropolis Algorithm
\end{itemize}

\section*{応用：ベイズ推測法}\label{ux5fdcux7528ux30d9ux30a4ux30baux63a8ux6e2cux6cd5}
\addcontentsline{toc}{section}{応用：ベイズ推測法}

\begin{itemize}
\tightlist
\item
  Bayesian Markov chain Monte Carlo method
\item
  STAN programming: \url{https://mc-stan.org/docs/reference-manual/mcmc.html}
\item
  Using R: \url{https://jtr13.github.io/cc20/a-basic-introduction-to-markov-chain-monte-carlo-method-in-r.html}
\end{itemize}

\chapter{Sufficient Statistics（十分統計量）}\label{sufficient-statisticsux5341ux5206ux7d71ux8a08ux91cf}

\section*{復習}\label{ux5fa9ux7fd2-4}
\addcontentsline{toc}{section}{復習}

\begin{itemize}
\tightlist
\item
  Exponential family
\item
  Statistic
\item
  Likelihood
\end{itemize}

\section*{Principles of data reduction}\label{principles-of-data-reduction}
\addcontentsline{toc}{section}{Principles of data reduction}

\begin{itemize}
\tightlist
\item
  Sufficiency principle
\item
  Likelihood principal
\item
  Formal likelihood Principal*

  \begin{itemize}
  \tightlist
  \item
    Formal sufficiency principle
  \item
    Conditional principle
  \item
    Formal likelihood principal
  \end{itemize}
\item
  Equivariance Principle*
\end{itemize}

\section{Sufficiency Principle}\label{sufficiency-principle}

\section{Sufficient Statistics}\label{sufficient-statistics}

\section{Ancillary Statistics}\label{ancillary-statistics}

\section{Complete Statistics}\label{complete-statistics}

\section{Basu' Theorem}\label{basu-theorem}

\section{Likelihood Principle}\label{likelihood-principle}

\section*{応用：ベイズ推測法}\label{ux5fdcux7528ux30d9ux30a4ux30baux63a8ux6e2cux6cd5-1}
\addcontentsline{toc}{section}{応用：ベイズ推測法}

\begin{itemize}
\tightlist
\item
  Bayesian Markov chain Monte Carlo method
\item
  STAN programming: \url{https://mc-stan.org/docs/reference-manual/mcmc.html}
\item
  Using R: \url{https://jtr13.github.io/cc20/a-basic-introduction-to-markov-chain-monte-carlo-method-in-r.html}
\end{itemize}

\end{document}
